\providecommand{\itescastudentname}{Martin Jonathan de la Cruz}
\providecommand{\itescastudentid}{Por confirmar}
\providecommand{\itescastudentemail}{Dato no proporcionado}
\providecommand{\itescafaculty}{Maestría en Gestión Administrativa}
\providecommand{\itescacoursename}{Seminario I}
\providecommand{\itescacoursecode}{MGA-SEMINARIO-I}
\providecommand{\itescadepartment}{Posgrado e Investigación}
\providecommand{\itescateachingfigure}{Dra. Carla Olimpya Zapuche Moreno}
\providecommand{\itescaprogramlevel}{Maestría}
\providecommand{\itescaprogramperiod}{Agosto--noviembre de 2026}
\providecommand{\itescamodality}{En línea}
\providecommand{\itescalocation}{Cajeme, Sonora}
\providecommand{\itescaasset}{ITESCA/assets-itesca/itesca-monograma.png}
\providecommand{\itescaofficialasset}{ITESCA/assets-itesca/web/logo-itesca-contorno.png}
\providecommand{\itescasecondaryasset}{ITESCA/assets-itesca/itesca-monograma.png}
\providecommand{\itescabibliography}{ITESCA/maestria-en-gestion-administrativa/seminario-i-mga/seminario-i}

\def\itescadocumenttitle {Mi ruta hacia la obtención del grado de Maestro en Gestión Administrativa}
\def\itescadocumentsubtitle {Actividad 1}
\def\itescadocumentsubject {Ruta académica y profesional de titulación}
\def\itescaactivityname {Actividad 1. Mi ruta hacia la titulación en la MGA}
\def\itescasessionname {Unidad 1}
\def\itescablockname {Seminario I}
\def\itescamainproduct {Cuadro de integración y ensayo reflexivo}

\documentclass[spanish,letterpaper,oneside]{article}

\def\documenttitle {\itescadocumenttitle}
\def\documentsubtitle {\itescadocumentsubtitle}
\def\documentsubject {\itescadocumentsubject}
\def\documentauthor {\itescastudentname}
\def\coursename {\itescacoursename}
\def\coursecode {\itescacoursecode}
\def\universityname {Instituto Tecnológico Superior de Cajeme}
\def\universityfaculty {\itescafaculty}
\def\universitydepartment {\itescadepartment}
\def\universitydepartmentimage {\itescaofficialasset}
\def\universitydepartmentimagecfg {width=5.2cm}
\def\universitylocation {\itescalocation}

\def\authortable {
  \begin{tabular}{ll}
    \textbf{Alumno:} & \itescastudentname \\
    No. de control: & \itescastudentid \\
    Carrera: & \itescafaculty \\
    Semestre: & \itescaprogramperiod \\
    Asignatura: & \itescacoursename \\
    Docente: & \itescateachingfigure \\
    Modalidad: & \itescamodality \\
    Fecha de entrega: & \today \\
    Lugar: & \universitylocation
  \end{tabular}
}

\input{template}
\usepackage{helvet}
\renewcommand{\familydefault}{\sfdefault}
\usepackage{array}
\usepackage{longtable}
\usepackage{ragged2e}
\usepackage{enumitem}
\usepackage{pdflscape}
\setcitestyle{authoryear,open={(},close={)}}
\setlist[itemize]{leftmargin=*,nosep}

\definecolor{itescaRedDark}{HTML}{8F1117}
\definecolor{itescaRed}{HTML}{D70F18}
\definecolor{itescaCharcoal}{HTML}{141414}
\definecolor{itescaSilver}{HTML}{D9D9D9}
\def\portraittitlecolor {itescaRedDark}

\begin{document}
\templatePortrait
\templatePagecfg
\onehalfspacing

\begin{abstractd}
\textbf{Objetivo:} Reflexionar sobre el perfil de egreso, los objetivos, las LGAC y los requisitos académicos de la Maestría en Gestión Administrativa para vincularlos con un proyecto de investigación y una ruta eficiente de titulación.

\textbf{Producto:} Cuadro de integración de doce ideas principales y ensayo reflexivo sobre la trayectoria personal, profesional y académica.
\end{abstractd}

\templateIndex
\templateFinalcfg

\section{Introducción}

Obtener el grado de Maestro en Gestión Administrativa exige mucho más que aprobar asignaturas: requiere convertir un interés profesional en un proyecto de tesis pertinente, original y útil para una organización. La ruta de titulación debe diseñarse desde el inicio porque el perfil de egreso, los objetivos del programa, las LGAC y los requisitos normativos forman un sistema: orientan qué problema investigar, con qué capacidades intervenirlo y qué evidencias producir para graduarse \citep{itescaMGA2026}. A partir de esta premisa, se presenta un cuadro de integración institucional y una reflexión personal enfocada en la mejora de la mercadotecnia y las ventas mediante decisiones sustentadas en datos.

\section{Ruta estratégica de formación y titulación}

\subsection{Cuadro de integración}

La síntesis siguiente recupera la información oficial de la MGA publicada por ITESCA y los requisitos establecidos en el numeral 3.7.2 de los Lineamientos para la Operación de los Estudios de Posgrado del TecNM, disponibles en el repositorio institucional de normatividad \citep{tecnmLineamientos2023,itescaNormatividad2026}. El cuadro se lee por filas: cada apartado reúne tres o más ideas institucionales y, en la última columna, su traducción a decisiones de formación, investigación y desempeño profesional. Los conceptos rectores son \textbf{perfil de egreso}, \textbf{objetivos del programa}, \textbf{LGAC}, \textbf{innovación organizacional}, \textbf{gestión financiera}, \textbf{tesis individual}, \textbf{originalidad} y \textbf{retribución social}.

\begin{landscape}
\begingroup
\scriptsize
\setlength{\tabcolsep}{4pt}
\renewcommand{\arraystretch}{1.18}
\begin{longtable}{@{}p{0.15\linewidth}p{0.48\linewidth}p{0.31\linewidth}@{}}
\caption{Integración de los componentes que orientan la trayectoria en la MGA}\label{tab:integracion-mga}\\
\toprule
\textbf{Apartado} & \textbf{Ideas principales} & \textbf{Relación personal y profesional} \\
\midrule
\endfirsthead
\toprule
\textbf{Apartado} & \textbf{Ideas principales} & \textbf{Relación personal y profesional} \\
\midrule
\endhead
\textbf{Perfil de egreso} &
\begin{itemize}
 \item Resolver problemas de gestión organizacional, calidad y competitividad mediante TIC.
 \item Analizar el entorno económico local, nacional e internacional considerando factores sociales, políticos y culturales.
 \item Intervenir procesos, estructuras y tareas, así como crear, dirigir o asesorar organizaciones con enfoque transdisciplinar, ético y sustentable.
\end{itemize} &
Estas capacidades se vinculan con mi propósito de transformar datos comerciales en decisiones sobre clientes, procesos de venta y competitividad. El componente ético evita que la innovación comercial se reduzca a vender más sin valorar sus efectos en consumidores, personal y entorno. \\
\midrule
\textbf{Objetivos generales y específicos} &
\begin{itemize}
 \item Formar líderes capaces de intervenir y solucionar problemas relevantes de empresas y sectores estratégicos regionales.
 \item Aplicar fundamentos administrativos, financieros y estadísticos con ética y responsabilidad social para elevar calidad y productividad.
 \item Gestionar recursos, innovación tecnológica, proyectos, cadena de suministro y calidad; además, formular diagnósticos, planes de negocio y programas de mejora con métodos de investigación.
\end{itemize} &
El programa fortalece mi transición de una práctica comercial basada en intuición hacia una gestión sustentada en evidencia. Quiero integrar análisis estadístico, diagnóstico organizacional y tecnologías digitales para proponer mejoras medibles en mercadotecnia y ventas. \\
\midrule
\textbf{LGAC y alcance} &
\begin{itemize}
 \item \emph{Gestión e Innovación de las Organizaciones}: estudia innovación abierta y capacidades para innovar en productos, procesos, organización, comercialización, tecnología y modelos de negocio.
 \item Esta línea analiza impactos económicos, sociales y ambientales, especialmente en micro, pequeñas y medianas empresas y en redes regionales de colaboración.
 \item \emph{Gestión económica y financiera de las organizaciones}: aborda planeación, organización, dirección y control de recursos, presupuestos, riesgos e información financiera para el crecimiento sostenible.
\end{itemize} &
La primera LGAC es la más cercana a mi interés porque permite investigar innovación comercial, experiencia del cliente y mejora del proceso de ventas. La segunda funcionará como apoyo para evaluar viabilidad, riesgo y retorno de las acciones propuestas. \\
\midrule
\textbf{Requisitos académicos para el grado} &
\begin{itemize}
 \item Acreditar el programa con promedio general mínimo de 80 y cumplir las actividades académicas.
 \item Desarrollar tesis individual congruente con una LGAC, con carta de originalidad y similitud total menor al 30\%; presentar carta de usuario sobre incidencia y satisfacción.
 \item Acreditar un segundo idioma, realizar al menos dos actividades de retribución social y generar un producto académico original derivado de la tesis.
 \item Obtener autorización de impresión de la DEPI y aprobar el examen de grado mediante la defensa de tesis.
\end{itemize} &
Asumo estos requisitos como entregables de un solo proyecto y no como trámites aislados. Desde el primer semestre organizaré un expediente de evidencias, calendario de idioma, retribución social, producto académico, revisiones de originalidad y validación con la organización usuaria. \\
\bottomrule
\end{longtable}
\endgroup
\end{landscape}

\subsection{Ensayo reflexivo: una ruta que comienza con el problema y termina con evidencia}

Mi proyecto personal y profesional se orienta a mejorar la gestión de mercadotecnia y ventas mediante información confiable, procesos medibles y decisiones centradas en el cliente. Esta aspiración coincide con el perfil de egreso de la MGA porque el programa no limita la administración a funciones operativas: exige resolver problemas, comprender el entorno e intervenir procesos y estructuras con apoyo de tecnologías digitales. Para mí, esa combinación significa avanzar de la experiencia práctica a una capacidad profesional para diagnosticar por qué una estrategia comercial funciona, reconocer sus límites y demostrar el efecto de una mejora.

Los objetivos de la maestría refuerzan esta dirección. El uso articulado de fundamentos administrativos, gestión financiera y estadística aplicada puede ayudarme a convertir datos de ventas, comportamiento del cliente y desempeño de campañas en evidencia para decidir. Asimismo, la ética, la responsabilidad social y la sustentabilidad ofrecen criterios indispensables: una intervención comercial debe ser rentable, pero también transparente, respetuosa de las personas y sostenible. Por ello, visualizo como proyecto inicial el diseño y evaluación de una estrategia de innovación comercial para una micro o pequeña empresa, con indicadores de conversión, retención, satisfacción y rentabilidad.

Desde mi perspectiva, la LGAC más cercana es \emph{Gestión e Innovación de las Organizaciones}. Su atención a innovaciones de comercialización, procesos y modelos de negocio permite estudiar la experiencia del cliente y la transformación digital sin aislarlas del sistema organizacional. La LGAC de \emph{Gestión económica y financiera de las organizaciones} complementaría el estudio al valorar presupuesto, riesgo y resultados. El portal identifica a la Mtra. Manuela Ruiz Castro como responsable de la MGA y al Dr. Luis Alberto Limón Valencia en la Subdirección de Posgrado e Investigación \citep{itescaMGA2026}; los considero primeros enlaces para solicitar el directorio del núcleo académico y elegir, con evidencia de trayectoria, a quien realmente trabaje la línea y el problema propuestos. A mi juicio, no sería riguroso atribuirles una especialidad que la página consultada no publica.

Para cumplir eficientemente los requisitos de titulación aplicaré una estrategia de gestión por hitos. Durante el primer semestre delimitaré problema, organización usuaria, LGAC, pregunta, indicadores y cronograma con el posible director. En paralelo iniciaré la acreditación del segundo idioma y programaré las dos actividades de retribución social para evitar que se acumulen al final. Cada asignatura deberá aportar un insumo reutilizable para la tesis: revisión teórica, instrumento, análisis o procedimiento \citep{tecnmLineamientos2023}.

También mantendré una matriz mensual con promedio, avances, acuerdos tutoriales, versiones y respaldos; usaré un gestor bibliográfico y revisaré la similitud desde los borradores, no únicamente antes de imprimir. Acordaré con la organización la evidencia necesaria para obtener la carta de usuario y planearé desde temprano un producto académico derivado. Finalmente, reservaré un periodo para dictamen, autorización de impresión y simulacros de defensa. Así, titularme deja de ser una meta distante y se convierte en una secuencia controlable de productos conectados con una intervención útil.

\clearpage
\section{Conclusiones}

Sostengo que la ruta hacia el grado depende de alinear desde el inicio el proyecto profesional, una LGAC y cada requisito normativo. Mi interés en mercadotecnia y ventas encuentra su mejor encuadre en la innovación organizacional, complementada por evaluación económica y financiera. Por consiguiente, la eficiencia no dependerá de acelerar trámites al final, sino de producir durante toda la maestría evidencia original, útil y verificable. Mantener un cronograma de hitos, vincular las asignaturas con la tesis y dialogar oportunamente con la coordinación y el núcleo académico permitirá que la investigación contribuya tanto a mi formación como a la organización participante. La responsabilidad sobre la exactitud, selección de fuentes y postura final permanece en el autor, quien empleó inteligencia artificial como apoyo de organización, revisión y edición.\footnote{Declaración de uso: se utilizó inteligencia artificial como apoyo para planificar, verificar consistencia y revisar la presentación; las fuentes institucionales fueron consultadas y las decisiones argumentativas fueron validadas por el autor.}

\clearpage
\bibliographystyle{apalike}
\bibliography{\itescabibliography}

\end{document}
